\section{Funktioner - The Basics}
En funktion er et matematisk udtryk, der beskriver sammenhængen mellem en eller flere variable. Simple funktioner opererer typisk med én uafhængig variabel $x$ og en afhængig variabel $f(x)$. Den afhængige variabel afhænger af $x$, hvilket ses i notationen $f(x)$.

Nogle eksempler på funktioner er:
\begin{align*}
    f(x) &= x \\
    g(x) &= 1 + x \\
    h(x) &= x^2 \\
    k(x) &= 5x + 1 - 7x^4 \\
    l(t) &= 3t + 1
\end{align*}

En funktion er et matematisk udtryk, der bruges til at beskrive en sammenhæng. Vi kan f.eks. beskrive prisen på en taxa-tur. Antag, at der er et fast gebyr på 100 DKK, og derefter koster det 10 DKK pr. kilometer. Funktionen kan skrives som:
\begin{equation*}
    T(x) = 100 + 10x
\end{equation*}
Her er $T(x)$ den samlede pris, og $x$ er antallet af kørte kilometer. Prisen afhænger altså af, hvor langt der køres.

\subsection{Definitionsmængde}
Definitionsmængden er den mængde af værdier, som den uafhængige variabel kan antage, uden at funktionen bliver udefineret.

Matematisk bestemmes definitionsmængden af, hvilke værdier der gør udtrykket gyldigt. I praksis kan man også overveje, hvilke værdier der giver mening i en konkret situation. I taxa-eksemplet giver negative værdier af $x$ ikke mening, da man ikke kan køre en negativ afstand.

Rent matematisk er definitionsmængden:
\begin{equation*}
    Dm(T) = ]-\infty, \infty[
\end{equation*}
Hvis vi kun ønsker realistiske værdier (ikke-negative afstande), kan vi skrive:
\begin{equation*}
    Dm(T) = [0, \infty[
\end{equation*}

Klammerne angiver, om endepunkter er med:
\begin{itemize}
    \item $[a,b]$ inkluderer både $a$ og $b$
    \item $]a,b[$ ekskluderer både $a$ og $b$
\end{itemize}
Uendelighed har altid åbne klammer: $]-\infty, \infty[$.

Et eksempel på en funktion, der ikke er defineret for alle $x$, er:
\begin{equation*}
    f(x) = \frac{1}{x}
\end{equation*}
Man kan ikke dividere med 0, så:
\begin{equation*}
    Dm(f) = ]-\infty, 0[ \cup ]0, \infty[
\end{equation*}
Her betyder $\cup$ union (forening af mængder).

Et andet eksempel:
\begin{equation*}
    g(x) = \frac{1}{x-3}
\end{equation*}
Her må $x \neq 3$, så:
\begin{equation*}
    Dm(g) = ]-\infty, 3[ \cup ]3, \infty[
\end{equation*}

\subsection{Værdimængde}
Værdimængden er de værdier, funktionen kan antage. Man kan forestille sig, at man indsætter alle værdier fra definitionsmængden i funktionen; de værdier, man får ud, udgør værdimængden.

Eksempel:
\begin{equation*}
    f(x) = 5x + 3 \qquad Dm(f) = ]-\infty, \infty[
\end{equation*}
Funktionen kan antage alle reelle værdier, så:
\begin{equation*}
    Vm(f) = ]-\infty, \infty[
\end{equation*}

Et andet eksempel:
\begin{equation*}
    g(x) = 2 + 3x^2 \qquad Dm(g) = ]-\infty, \infty[
\end{equation*}
Da $x^2 \geq 0$, vil $g(x) \geq 2$. Minimum opnås for $x=0$:
\begin{equation*}
    g(0) = 2
\end{equation*}
Funktionen kan vokse uendeligt, så:
\begin{equation*}
    Vm(g) = [2, \infty[
\end{equation*}

\subsection{Identitetsfunktionen}
Identitetsfunktionen er den simple funktion
\begin{equation*}
    f(x) = x
\end{equation*}
Det vil sige, at funktionen returnerer den samme værdi, som den får som input.

Generelt gælder, at enhver funktion på formen $f(x) = x$ (eller f.eks. $g(x) = x$) kaldes en identitetsfunktion.

Med andre ord: Den afhængige variabel antager præcis samme værdi som den uafhængige variabel.